% Generated from tex/chapters/ch04/07_実務上の考慮.tex for the SWEBOK structure tree
\subsubsection{コーディング}

コーディングでは、ソースコードを理解しやすく、変更しやすく、検証しやすく書くことが重要である。単に動作するコードを書くだけでは不十分である。コードは、将来の読者に意図を伝える文書でもある。

構築におけるコーディングの主な考慮事項は次のとおりである。

\begin{itemize}
  \item 命名規則、字下げ、レイアウトにより、理解しやすいソースコードを作る。
  \item クラス、列挙型、変数、名前付き定数などを適切に使う。
  \item 条件分岐、ループ、例外処理などの制御構造を過度に複雑にしない。
  \item 予期できるエラーと例外的なエラーを区別して扱う。
  \item バッファオーバーフロー、配列範囲外参照、入力検証不足など、コードレベルのセキュリティ問題を防ぐ。
  \item スレッド、データベースロック、ファイル、ネットワーク接続などの資源を適切に扱う。
  \item コードを関数、クラス、パッケージなどに整理する。
  \item 必要な文書化とコメントを行う。
  \item 測定に基づいて性能を調整する。
\end{itemize}

\begin{keybox}{コメントの考え方}
コメントは、コードの日本語訳ではない。よいコメントは、コードだけでは読み取りにくい理由、前提、制約、例外、意図を説明する。名前と構造で伝えられることは、名前と構造で伝える方がよい。
\end{keybox}
